\documentclass{chsmc}
\chsmcsetup{
  year = 2021,
  part = 1,
  type = B,
  show-answers = false,
}
\begin{document}

\section{填空题：本大题共 8 小题，每小题 8 分，共 64 分.}

\begin{enumerate}
  \item 等差数列 $\{a_n\}$ 的公差 $d \neq 0$，且 $a_{2021} = a_{20} + a_{21}$，
    则 $\dfrac{a_1}{d}$ 的值为 \blankline
  \item 设 $m$ 为实数，复数 $z_1 = 1 + 2\ii$，$z_2 = m + 3\ii$
    （这里 $\ii$ 为虚数单位），若 $z_1\cdot\bar{z}_2$ 为纯虚数，
    则 $|z_1 + z_2|$ 的值为 \blankline
  \item 当 $\dfrac{\pi}{4} \leqslant x \leqslant \dfrac{\pi}{2}$ 时，
    $y = \sin^2 x + \sqrt{3}\sin x\cos x$ 的取值范围是 \blankline
  \item 设函数 $f(x)$ 的定义域为 $D = (-\infty,0)\cup(0,+\infty)$，
    且对任意 $x \in D$，均有
    $f(x) = \dfrac{f(1)\cdot x^2 + f(2)\cdot x - 1}{x}$，
    则 $f(x)$ 的所有零点之和为 \blankline
  \item 设 $a, b, c > 1$，满足
    $(a^2 b)^{\log_a c} = a\cdot(ac)^{\log_a b}$，
    则 $\log_c(ab)$ 的值为 \blankline
  \item 在 $\triangle ABC$ 中，$AB = 1$，$AC = 2$，$\cos B = 2\sin C$，
    则 $BC$ 的长为 \blankline
  \item 在平面直角坐标系 $xOy$ 中，已知抛物线
    $y = ax^2 - 3x + 3$ $(a \neq 0)$ 的图象与抛物线 $y^2 = 2px$ $(p > 0)$
    的图象关于直线 $y = x + m$ 对称，则实数 $a$, $p$, $m$ 的乘积为 \blankline
  \item 设 $a_1, a_2, \cdots, a_{10}$ 是 $1, 2, \cdots, 10$ 的一个随机排列，
    则在 $a_1 a_2, a_2 a_3, \cdots, a_9 a_{10}$ 这 $9$ 个数中既出现 $9$
    又出现 $12$ 的概率为 \blankline
\end{enumerate}

\section{解答题：本大题共 3 个小题，满分 56 分. 解答应写出文字说明、证明过程或演算步骤.}

\begin{enumerate}[resume]
  \item (本题满分 16 分) 数列 $\{a_n\}$ 满足：$a_1 = a_2 = a_3 = 1$.
    令 $b_n = a_n + a_{n+1} + a_{n+2}$ $(n \in \mathbb{N}^*)$.
    若 $\{b_n\}$ 是公比为 $3$ 的等比数列，求 $a_{100}$ 的值.
  \item (本题满分 20 分) 在平面直角坐标系中，函数
    $y = \dfrac{1}{|x|}$ 的图象为 $\Gamma$. 设 $\Gamma$ 上的两点 $P$, $Q$ 满足：
    $P$ 在第一象限，$Q$ 在第二象限，且直线 $PQ$ 与 $\Gamma$ 位于第二象限的部分相切于点 $Q$.
    求 $|PQ|$ 的最小值.
  \item (本题满分 20 分) 在正 $n$ $(n \geqslant 3)$ 棱锥
    $P\text{-}A_1 A_2 \cdots A_n$ 中，$O$ 为底面正 $n$ 边形
    $A_1 A_2 \cdots A_n$ 的中心，$B$ 为棱 $A_1 A_n$ 的中点.
    \begin{enumerate}[(1)]
      \item 证明：$PO^2\sin^2\dfrac{\pi}{n} + PA_1^2\cos^2\dfrac{\pi}{n} = PB^2$；
      \item 设正 $n$ 棱锥 $P\text{-}A_1 A_2 \cdots A_n$ 的侧棱与底面所成的角为 $\alpha$，
        侧面与底面所成的角为 $\beta$，试确定
        $\dfrac{1}{n}\displaystyle\sum_{i=1}^{n}\cos\angle A_i PB$
        与 $\sin\alpha\cdot\sin\beta$ 的大小关系，并予以证明.
    \end{enumerate}
\end{enumerate}

\end{document}